\section{Research Roadmap}
\label{sec:roadmap}

The BP interpretation of transformers is an active empirical program, not
a finished doctrine. The open research directions divide into three tiers
by tractability.

\subsection*{Tier 1 --- Most Directly Testable}

\textbf{Recover factor graphs from SAEs.} For each boolean-AND head, project
the Q/K matrix onto the SAE feature basis and check if the dominant
component corresponds to a monosemantic feature. Check if the V circuit
routes from a belief dimension to a scratch slot.

\textbf{Test layer-wise belief convergence.} Train a tiny transformer on a
synthetic Bayes net with known ground-truth posteriors. Measure how closely
the residual stream at each layer matches the BP belief after that many
rounds of message passing.

\textbf{Measure calibration by layer depth.} Compare output calibration of
models with different depths on multi-step reasoning tasks. If deeper models
are better calibrated on longer chains, that supports the
layers-as-BP-rounds account.

\subsection*{Tier 2 --- Requires New Infrastructure}

Complete SAE coverage of a production model. Build and benchmark a grounded
transformer + QBBN hybrid. Develop a calibration benchmark for BP
consistency that measures layer-wise belief stability.

\subsection*{Tier 3 --- Longer-Term}

Learning grounded factor graphs from text via expectation maximization.
Formal account of in-context learning as BP. Multi-agent BP convergence
conditions.

\subsection*{What Would Most Advance the Program}

In order of expected impact: (1) the SAE factor graph recovery experiment,
which closes the loop between the formal theorem and trained models; (2)
layer-wise convergence measurement, which tests the layers-as-rounds claim
directly; and (3) the grounded hybrid system benchmark, which demonstrates
the practical value of grounding.